% §17.2 The Bernoulli Framework as Error Model
%
% Monograph register: cipher-map rates ARE Bernoulli rates; per-query
% channel matrices; chained composition is matrix multiplication
% (thm:cipher-composition); perfect-filter geometry carries over;
% two-stage encrypted-classifier worked example. Tightened. Notes the
% trapdoor set as a Bernoulli<T> instance (salvaged ex-10 point, inline).
% Cross-refs sec:noisy-bit, sec:bernoulli-bool-formal, sec:kronecker,
% thm:kronecker, ax:element-indep, sec:perfect-filter-line,
% def:perfect-filter-line, sec:n-ary-channels, def:n-ary-channel,
% thm:n-ary-kronecker, def:cipher-map, sec:bernoulli-of-composite.
% DEFINED: sec:bernoulli-error-model, thm:cipher-composition.

\section{The Bernoulli Framework as Error Model}
\label{sec:bernoulli-error-model}

\Cref{def:cipher-map} stated the abstraction in terms of a
false-positive rate $\fprate$ and a false-negative rate $\fnrate$. That
notation was not cosmetic: the rates are the same objects chapters 1 and
4 introduced for the Bernoulli atom, and a chained cipher-map composition
obeys the same laws as a chained noisy-classifier composition. This
section makes the equivalence operational, so the algebraic identities
the framework proved earlier carry over to cipher maps unchanged.

\subsection*{Cipher-map rates are Bernoulli rates}

A Bernoulli predicate has per-query behavior summarized by $\fprate =
\Prob{\tilde A(x) = 1 \Given x \notin A}$ and $\fnrate = \Prob{\tilde
A(x) = 0 \Given x \in A}$, and \Cref{sec:bernoulli-bool-formal} showed
that the two rates with the input prevalence specify the joint
distribution of any finite query sequence. A cipher map $\tilde f$ for a
Boolean predicate $f$ has the same per-query behavior. The client's only
observable is the decrypted bit $\mathsf{Dec}(\tilde f(\mathsf{Enc}(x)))
\in \{0, 1\}$, and that bit is Bernoulli conditional on the plaintext
class:
\[
  \fprate = \Prob{\mathsf{Dec}(\tilde f(\mathsf{Enc}(x))) = 1 \Given f(x) = 0},
  \quad
  \fnrate = \Prob{\mathsf{Dec}(\tilde f(\mathsf{Enc}(x))) = 0 \Given f(x) = 1},
\]
the probability taken over $\tilde f$'s internal randomness (hash
choice, ciphertext noise, key schedule). Encryption and decryption drop
out because $\mathsf{Dec} \circ \mathsf{Enc}$ is the identity on
correctly decrypted outputs, leaving a Bernoulli atom with the cipher
map's error profile as its $(\fprate, \fnrate)$. Every theorem of
chapters 1 through 15 that took $(\fprate, \fnrate)$ as input applies
when the rates come from a cipher map; the framework is indifferent to
where the error originates and tracks only how it propagates. A trapdoor
set is in this sense a $\Bernoulli{T}$ over its plaintext type $T$
(\Cref{sec:bernoulli-of-composite}), the encryption layer changing the
representation but not the Bernoulli-square point.

\subsection*{Per-query channel matrices}

The two-rate summary is convenient but limited. A Boolean cipher map is
fully captured by a $2 \times 2$ channel matrix
\[
  Q \;=\;
  \begin{pmatrix}
    1 - \fprate & \fprate \\
    \fnrate & 1 - \fnrate
  \end{pmatrix},
\]
rows indexed by the true plaintext class, columns by the observed
output; the Bloom-based map has $\fnrate = 0$, a zero in the lower-left.
For non-Boolean cipher maps (encrypted classifiers, ranked retrieval)
the $2 \times 2$ matrix generalizes to the $K \times L$ channel matrix
of \Cref{def:n-ary-channel}, with $Q_{k,\ell} = \Prob{\mathsf{Dec}(\tilde
f(\mathsf{Enc}(x))) = \ell \Given f(x) = k}$. The channel-matrix
representation captures the full per-query behavior in one object that
composes by multiplication, and it makes contact with the
information-theoretic machinery of chapter 8 that
\Cref{sec:maximizing-confidentiality} needs to quantify confidentiality.

\subsection*{Chained composition is matrix multiplication}

Suppose the client computes $h = f \circ g$ on encrypted data, $X
\xrightarrow{g} Y \xrightarrow{f} Z$: a cipher map $\tilde g$ returns an
intermediate $\mathsf{Enc}(Y)$, then $\tilde f$ computes on it and
returns $\mathsf{Enc}(Z)$. The error compounds. If $\tilde g$ produces a
wrong $y' \ne g(x)$, then $\tilde f$ is asked to compute $f(y')$, and
even an errorless $\tilde f$ returns $f(y') \ne f(g(x))$: the stage-1
error propagates through stage 2. This is the situation the chapter 4
Kronecker theorem (\Cref{thm:kronecker}) and its n-ary generalization
(\Cref{thm:n-ary-kronecker}) were built to analyze, and the
factorization is matrix multiplication.

\begin{theorem}[Cipher-map composition]
\label{thm:cipher-composition}
Let $\tilde g$ and $\tilde f$ be cipher maps realizing $g : X \to Y$ and
$f : Y \to Z$ with per-query channel matrices $Q_g \in [0,1]^{|X| \times
|Y|}$ and $Q_f \in [0,1]^{|Y| \times |Z|}$. If \Cref{ax:element-indep}
(Axiom 1) holds, so the two stages' errors are mutually independent
given the inputs, then the composed cipher map $\tilde f \circ \tilde g$
realizes $h = f \circ g$ with channel matrix
\[
  Q_h \;=\; Q_g \, Q_f,
\]
ordinary matrix multiplication, $Q_h \in [0,1]^{|X| \times |Z|}$.
\end{theorem}

The proof reuses the chapter 4 argument: $(Q_g Q_f)_{x,z}$ sums over
intermediate $y$ the probability that stage 1 produces $y$ times the
probability that stage 2 then produces $z$, the conditional independence
being Axiom 1 applied to the two-stage sequence. For Boolean maps the
theorem is a $2 \times 2$ product; the same result governs any finite
chain by induction, the composed matrix being the product of the
per-stage matrices. A consequence: chained cipher-map error grows
predictably, the composed Boolean rates bounded by chapter 3's
$1 - \prod_i (1 - \eta_i)$ on each error type. This is Property 4 of the
cipher-map design in the trapdoor-computing literature, the framework's
composition law and the cryptographic side's requirement being the same
property.

\subsection*{The geometry carries over}

The perfect-filter line of \Cref{def:perfect-filter-line} classifies
approximate sets by their $(\fprate, \fnrate)$ position, and the same
classification applies to cipher maps. A cipher map with $\fnrate = 0$ is
\emph{safe}: it never misses a real positive, the client losing no real
results, at the cost of $\fprate > 0$. The Bloom-based map sits on this
line, as does any cipher map built by encrypting a perfect-filter-line
approximate set. A cipher map with $\fprate = 0$ sits on the opposite
axis, \emph{exact when positive} but missing some real positives;
it appears in randomized-encoding maps that occasionally reject valid
plaintexts. The classification is geometric and applies to cipher maps as
a special case of approximate functions.

\subsection*{Worked example: two-stage encrypted classifier}

A representative composition appears in encrypted retrieval. The client
retrieves and ranks documents matching an encrypted query in two stages.
\emph{Stage 1, encrypted membership}: the server holds an encrypted Bloom
filter on a candidate set $A$, returns an encrypted membership bit at
$\fprate_1 = 0.01$, $\fnrate_1 = 0$, so
\[
  Q_g = \begin{pmatrix} 0.99 & 0.01 \\ 0 & 1 \end{pmatrix}.
\]
\emph{Stage 2, exact ranking on positives}: an FHE-based ranking
introduces no error, $Q_f = I$. By \Cref{thm:cipher-composition}, $Q_h =
Q_g I = Q_g$: the composed map inherits $\fprate = 0.01$, $\fnrate = 0$.
Had stage 2 carried its own FPR-only error $\fprate_2 = 0.005$, the
product would give composed $\fprate = 0.99 \cdot 0.005 + 0.01 \cdot 1
\approx 0.015$, matching $1 - (1 - 0.01)(1 - 0.005) \approx 0.015$, with
$\fnrate$ staying zero because both stages are one-sided. The framework's
composition law applies to cipher maps with no modification. The next
section unpacks the operational consequences: when $\fprate, \fnrate$ are
also confidentiality parameters, the same matrix algebra governs how
confidentiality accumulates across composed queries.
